\chude{Thử}
\man{2}{Thông hiểu}
\begin{exsx}
Tìm $x$, biết $\dfrac{2}{3}x-\dfrac{1}{2}=\dfrac{5}{6}$. Sắp xếp các bước thành lời giải đúng.
\buoc{$\dfrac{2}{3}x=\dfrac{5}{6}+\dfrac{1}{2}$}
\buoc{$\dfrac{2}{3}x=\dfrac{4}{3}$}
\buoc{$x=\dfrac{4}{3}:\dfrac{2}{3}$}
\buoc{$x=2$}
\nhieu{$\dfrac{2}{3}x=\dfrac{5}{6}-\dfrac{1}{2}$}{Chuyển $-\dfrac12$ sang vế phải phải đổi dấu thành $+\dfrac12$.}
\nhieu{$x=\dfrac{4}{3}\cdot\dfrac{2}{3}$}{Tìm thừa số chưa biết: lấy tích chia cho thừa số đã biết, không nhân.}
\loigiai{Vậy $x=2$.}
\end{exsx}
\begin{exsx}
Tính $A=\left(\dfrac12+\dfrac13\right)\cdot\left(\dfrac12-\dfrac13\right)$.
\buoc[1]{$\dfrac12+\dfrac13=\dfrac56$}
\buoc[1]{$\dfrac12-\dfrac13=\dfrac16$}
\buoc{$A=\dfrac56\cdot\dfrac16$}
\buoc{$A=\dfrac{5}{36}$}
\nhieu{$\dfrac12+\dfrac13=\dfrac25$}{Không cộng tử với tử, mẫu với mẫu. Phải quy đồng: $\dfrac36+\dfrac26=\dfrac56$.}
\nhieu{$A=\dfrac56+\dfrac16$}{Giữa hai dấu ngoặc là phép nhân, không phải phép cộng.}
\loigiai{Hai bước tính trong ngoặc làm trước bước nào cũng đúng.}
\end{exsx}
\begin{exsd}
Cho $\triangle ABC$ có $AB=AC$, $M$ là trung điểm của $BC$. Chứng minh $AM\perp BC$ (góc $\angle AMB$).
\nut{1}{$AM\perp BC$}{hai góc kề bù bằng nhau nên mỗi góc bằng $90^\circ$}
\nut{2}[1]{$\widehat{AMB}=\widehat{AMC}$}{hai góc tương ứng}
\nut{3}[2]{$\triangle AMB=\triangle AMC$}{c.c.c}
\nut{4}[3]{$AB=AC$}{giả thiết}
\nut{5}[3]{$MB=MC$}{$M$ là trung điểm $BC$}
\nut{6}[3]{$AM$ là cạnh chung}{}
\an{3}
\anly{2,3}
\nhieu{$\triangle AMB=\triangle CMA$}{Viết sai thứ tự đỉnh tương ứng.}
\nhieu{$\triangle ABC=\triangle AMC$}{Hai tam giác này không bằng nhau.}
\nhieuly{hai cạnh tương ứng}{Ô bên trên là hai góc bằng nhau.}
\nhieuly{c.g.c}{Chưa có góc xen giữa.}
\end{exsd}
\begin{exsd}
Cho $\triangle ABC$ có $AB=AC$, $M$ là trung điểm của $BC$. Lấy $D$ trên cạnh $AB$, $E$ trên cạnh $AC$ sao cho $BD=CE$. Chứng minh $MD=ME$.
\nut{1}{$MD=ME$}{hai cạnh tương ứng}
\nut{2}[1]{$\triangle MBD=\triangle MCE$}{c.g.c}
\nut{3}[2]{$BD=CE$}{giả thiết}
\nut{4}[2]{$\widehat{DBM}=\widehat{ECM}$}{hai góc tương ứng}
\nut{5}[2]{$MB=MC$}{$M$ là trung điểm $BC$}
\nut{6}[4]{$\triangle ABM=\triangle ACM$}{c.c.c}
\nut{7}[6]{$AB=AC$}{giả thiết}
\nut{8}[6]{$MB=MC$}{$M$ là trung điểm $BC$}
\nut{9}[6]{$AM$ là cạnh chung}{}
\an{2,6}
\anly{2,4,6}
\nhieu{$\triangle MBD=\triangle MEC$}{Viết sai thứ tự đỉnh tương ứng.}
\nhieu{$\widehat{BMD}=\widehat{CME}$}{Chỉ suy ra được sau khi đã có hai tam giác bằng nhau.}
\nhieuly{c.c.c}{Chưa có $MD=ME$.}
\nhieuly{g.c.g}{Đó là c.g.c.}
\loigiai{Có thêm một nhánh.}
\end{exsd}
\man{3}{x}
\begin{exsx}
%@ a = ngaunhien(2, 9)
Tính $2x={{2*a}}$.
\buoc{$x={{2*a}}:2$}
\buoc{$x={{a}}$}
\nhieu{$x={{4*a}}$}{Nhân thay vì chia.}
\end{exsx}
